% Extracto editor-ready del contenido exclusivo de masas/16_hidrodinamica_y_gauge.tex.
% Las líneas 1--14 están publicadas byte-exactamente desde
% 74_rev2_realizaciones_i.tex:544--557. El propietario íntegro permanece
% preservado; este extracto conserva la continuación exclusiva.

\subsection{Comparación de las realizaciones nulas electromagnética y cromodinámica}
El fotón y los ocho gluones pertenecen a cartas nulas distintas. Para el fotón,

\[
m_\gamma=0
\]
en la rama electromagnética. Para los gluones,

\[
m_g=0
\]
en la representación adjunta de color. No se obtiene la nulidad tomando el límite de un carácter positivo; el dominio se bifurca antes. La realización fotónica exige Maxwell, invariancia de calibre y dos polarizaciones transversales. La realización gluónica exige la conexión de color, sus ocho generadores y la restricción al espacio físico de calibre.

La propagación bidireccional puede escribirse

\[
c_\pm=\frac{c}{1\pm\epsilon},
\qquad
\epsilon=q_-^{90}-q_+^{90},
\]
con media armónica

\[
\frac{2}{1/c_++1/c_-}=c.
\]
La velocidad bidireccional conserva \(c\); la orientación unidireccional conserva el sesgo. Una brecha adicional del registro genealógico puede realizarse mediante Proca--Stückelberg, pero no debe identificarse automáticamente con \(\epsilon\) ni con el desdoblamiento de la acción.

\Needspace{7\baselineskip}
